\bookchapter{บทที่ 7: วิเคราะห์ฟังก์ชันล้ำสมัยอย่างครบถ้วน}

\chapterintro{เทคโนโลยีก้าวหน้าไม่หยุดนิ่ง ฟังก์ชันของมือถือก็พัฒนาไปเรื่อย ๆ ทั้งการสื่อสารผ่านดาวเทียม การสื่อสารไร้เครือข่าย การตรวจจับอุบัติเหตุรถชน และเทคโนโลยี StarFlash ฟังก์ชันพวกนี้ฟังดูเหมือนนิยายวิทยาศาสตร์ แต่จริง ๆ ได้เข้ามาอยู่ในชีวิตเราแล้ว มาทำความรู้จักเทคโนโลยีล้ำหน้าเหล่านี้กัน}

\booksection{หนึ่ง การสื่อสารผ่านดาวเทียม: ไม่มีสัญญาณก็ขอความช่วยเหลือได้}

การสื่อสารผ่านดาวเทียมคือก้าวที่ยิ่งใหญ่ของวงการมือถือช่วงหลายปีที่ผ่านมา เวลาคุณอยู่กลางป่าลึกกลางดอน มือถือไม่มีสัญญาณเลยสักนิด การสื่อสารผ่านดาวเทียมอาจเป็นทางรอดเพียงทางเดียวของคุณก็ได้

\booksubsection{หลักการทำงาน:}

มือถือจะเชื่อมต่อตรงกับดาวเทียมสื่อสารที่โคจรอยู่บนฟ้า แล้วส่งข้อความขอความช่วยเหลือออกไป ปัจจุบันรองรับหลัก ๆ คือการส่งข้อความและตำแหน่ง ยังโทรออกไม่ได้

\booksubsection{รุ่นที่รองรับ:}

\begin{itemize}
  \item \textbf{Apple}:iPhone 14 ซีรีส์ขึ้นไป รองรับการส่ง SOS ฉุกเฉินผ่านดาวเทียม ในขณะที่ iPhone 17 ซีรีส์รองรับการส่ง iMessage และ SMS ผ่านดาวเทียมเพิ่มเติม พร้อมขยายการครอบคลุมไปยังพื้นที่เพิ่มเติม
  \item \textbf{แอปเปิล}: เริ่มตั้งแต่ Mate 50 ซีรีส์ที่รองรับข้อความผ่านดาวเทียม BeiDou (ข้อความล้วน) Mate 60 Pro ซีรีส์รองรับการโทรผ่านดาวเทียม Tiantong ได้จริง และ Mate 70 ซีรีส์ก็เก่งขึ้นไปอีกขั้น
  \item \textbf{Xiaomi/OPPO/vivo}: ตั้งแต่ปี 2024 เป็นต้นมา รุ่นฟลากชิปต่าง ๆ ก็เริ่มรองรับข้อความผ่านดาวเทียม BeiDou กันแล้ว
  \item \textbf{แนวโน้มพัฒนาการ}: การสื่อสารผ่านดาวเทียมกำลังขยายลงจากฟังก์ชันเฉพาะรุ่นฟลากชิปสู่รุ่นทั่วไป อนาคตรุ่นกลางก็น่าจะรองรับ ระบบ BeiDou ของจีนรองรับบนมือถือสัญชาติจีนกว้างขวางที่สุด
\end{itemize}

\booksubsection{สถานการณ์ที่ใช้:}

\begin{itemize}
  \item : ผจญภัยกลางธรรมชาติ ปีนเขาแล้วประสบภัย
  \item : ประสบอันตรายกลางทะเล
  \item : ภัยธรรมชาติทำให้การสื่อสารภาคพื้นดินล่ม
\end{itemize}

\booksubsection{วิธีใช้งาน:}

\begin{enumerate}
  \item : อยู่ในที่โล่งที่มองเห็นท้องฟ้า
  \item : เปิดฟังก์ชันการสื่อสารผ่านดาวเทียมของมือถือ (ปกติอยู่ในหน้าโทรฉุกเฉิน)
  \item : ทำตามขั้นตอนบนหน้าจอ แล้วเล็งไปทางดาวเทียม
  \item : ส่งข้อความขอความช่วยเหลือและตำแหน่งที่ตั้งไว้ล่วงหน้า
\end{enumerate}

\begin{tcbinfo}{เกร็ดความรู้}
การสื่อสารผ่านดาวเทียมต้องอาศัยเสาอากาศฮาร์ดแวร์เฉพาะในตัวมือถือ ไม่ใช่ทุกรุ่นจะรองรับ อีกอย่าง ปกติต้องสมัครสมาชิกแบบเสียเงิน แต่ฟังก์ชันขอความช่วยเหลือฉุกเฉินมักใช้ฟรี

\end{tcbinfo}

\booksection{สอง การสื่อสารไร้เครือข่าย: ทางใหม่ของการคุยระยะใกล้}

“การสื่อสารไร้เครือข่าย” คือเทคโนโลยีที่ให้มือถือสองเครื่องคุยกันได้ตรง ๆ ในสถานการณ์ที่ไม่มีทั้งเน็ตมือถือและ Wi-Fi

\booksubsection{หลักการทางเทคนิค:}

: ใช้บลูทูธหรือคลื่นวิทยุเฉพาะทาง สร้างการเชื่อมต่อแบบจุดต่อจุดในระยะใกล้ (ปกติไม่กี่ร้อยเมตรถึงไม่กี่กิโลเมตร)

\booksubsection{สถานการณ์ที่ใช้:}

\begin{itemize}
  \item \textbf{กิจกรรมกลางแจ้ง}: ตอนเดินป่า ตั้งแคมป์ ใช้คุยกันระหว่างเพื่อนร่วมทีม
  \item \textbf{การกู้ภัยภัยพิบัติ}: ช่วงแผ่นดินไหว น้ำท่วม เมื่อการสื่อสารล่ม
  \item \textbf{กิจกรรมคนชุกชุม}: คอนเสิร์ต งานกีฬา ช่วงที่คนแน่นจนสัญญาณแย่
\end{itemize}

\booksubsection{แพลตฟอร์มที่รองรับ:}

\begin{itemize}
  \item \textbf{Apple}: ฟังก์ชัน "NameDrop" และ "หาเพื่อนแบบออฟไลน์" ของ iOS 17
  \item \textbf{Huawei}: แอป ChangLian รองรับการโทรแบบไร้เครือข่าย
  \item \textbf{Xiaomi}: ฟังก์ชันบางส่วนของศูนย์แชร์อัจฉริยะ Xiaomi (Mi Smart Hub)
\end{itemize}

\booksection{สาม การตรวจจับอุบัติเหตุรถชน: ช่วยชีวิตได้ในวินาทีสำคัญ}

การตรวจจับอุบัติเหตุรถชนคือฟังก์ชันที่ใช้เซ็นเซอร์ของมือถือตรวจจับการชนรุนแรง แล้วโทรเรียกฉุกเฉินให้อัตโนมัติ

\booksubsection{หลักการทำงาน:}

มือถือประเมินจากเซ็นเซอร์หลายตัวพร้อมกัน ทั้งเซ็นเซอร์วัดแรงเร่ง ไจโรสโคป บารอมิเตอร์ และไมโครโฟน ว่าเกิดการชนรุนแรงขึ้นจริงหรือไม่ พอตรวจจับอุบัติเหตุได้ มือถือจะ:

\begin{enumerate}
  \item : ขึ้นหน้าจอนับเวลาถอยหลัง ถามว่าคุณสบายดีไหม
  \item : ถ้าหมดเวลานับถอยหลังแล้วคุณยังไม่ตอบสนอง จะโทรเรียกฉุกเฉินอัตโนมัติ
  \item : ส่งตำแหน่งของคุณให้ผู้ติดต่อฉุกเฉิน
\end{enumerate}

\booksubsection{รุ่นที่รองรับ:}

\begin{itemize}
  \item \textbf{Apple}: iPhone 14 ซีรีส์ขึ้นไป
  \item \textbf{Android}: รุ่นเรือธงบางรุ่นเริ่มรองรับเพิ่มขึ้นเรื่อย ๆ
\end{itemize}

\booksubsection{ข้อควรระวัง:}

\begin{itemize}
  \item : ฟังก์ชันต้องเปิดใช้งานก่อนถึงจะใช้ได้
  \item : อาจแจ้งเตือนผิด (อย่างตอนนั่งรถไฟเหาะหรือออกกำลังกายหักโหม) ทันเวลากดยกเลิกก็พอ
  \item : อาศัยเซ็นเซอร์และอัลกอริทึม ไม่ใช่ 100\% เชื่อถือได้ แต่ก็เป็นความปลอดภัยอีกชั้นหนึ่ง
\end{itemize}

\booksection{สี่ เทคโนโลยี StarFlash (SparkLink): การเชื่อมต่อไร้สายยุคถัดไป}

StarFlash (SparkLink) คือเทคโนโลยีสื่อสารไร้สายรุ่นใหม่ที่ Huawei นำทีมวิจัย ตั้งเป้ามาแทนบลูทูธและ Wi-Fi ในบางสถานการณ์ ในตัวชี้วัดสำคัญ ๆ ทัดเทียมหรือเกินโซลูชันที่มีอยู่: อัตราส่งข้อมูลได้ถึง 900Mbps (ไกลกว่าบลูทูธที่ 50Mbps) ความหน่วงต่ำถึง 20 $\mu$ s (บลูทูธประมาณ 15ms) และเครือข่ายเดียวเชื่อมอุปกรณ์พร้อมกันได้ 4096 ชิ้น (บลูทูธปกติไม่เกิน 8 ชิ้น)

\booksubsection{จุดเด่นทางเทคนิค:}

\begin{itemize}
  \item \textbf{ความหน่วงต่ำมาก}: ความหน่วงปลายทางถึงปลายทางแค่ 20 ไมโครวินาที (0.02ms) ต่ำกว่าบลูทูธถึงสองเท่าตัว เวลาเล่นเกมหรือฟังเพลงแทบไม่รู้สึกถึงความหน่วง
  \item \textbf{อัตราความเร็วสูง}: อัตราส่งข้อมูลสูงสุดต่อลิงก์ 900Mbps เป็นสิบเท่าของเสียงบลูทูธ ส่งเสียง lossless หรือแม้แต่วิดีโอความคมชัดสูงได้
  \item \textbf{ใช้พลังงานต่ำ}: ไฟสแตนด์บายเทียบเท่าบลูทูธ แต่ส่งข้อมูลได้มีประสิทธิภาพกว่า ใช้พลังงานต่อข้อมูลจำนวนเท่ากันน้อยกว่า
  \item \textbf{ครอบคลุมกว้าง + เชื่อมพร้อมกันได้เยอะ}: ระยะเชื่อมต่อไกลกว่าบลูทูธตั้ง 50\% ขึ้นไป เครือข่ายเดียวรองรับอุปกรณ์ออนไลน์พร้อมกันได้เป็นพัน ๆ ชิ้น
\end{itemize}

\booksubsection{สถานการณ์ที่ใช้:}

\begin{itemize}
  \item \textbf{หูฟังไร้สาย}: ความหน่วงต่ำลง เล่นเกมแล้วเสียงกับภาพเคลื่อนไหวตรงกัน
  \item \textbf{สมาร์ทวอทช์}: เชื่อมต่อเสถียรขึ้น
  \item \textbf{สมาร์ทโฮม}: เชื่อมอุปกรณ์พร้อมกันได้มากขึ้น
  \item \textbf{เชื่อมต่อในรถยนต์}: มือถือกับรถเชื่อมกันลื่นไหลขึ้น
\end{itemize}

\booksubsection{อุปกรณ์ที่รองรับ:}

ปัจจุบันเป็นอุปกรณ์ระบบนิเวศ Huawei เกือบทั้งหมด เช่น Huawei Mate 60 ซีรีส์ หูฟังและสมาร์ทวอทช์ Huawei บางรุ่น อนาคตจะมีแบรนด์อื่นรองรับเพิ่ม

\textbf{มุมมองเทคโนโลยี}\par

StarFlash เป็นตัวแทนของอนาคตการสื่อสารไร้สาย ยิ่งมีผู้ผลิตรายอื่นเข้าร่วมมากขึ้น เรายิ่งจะได้เห็นอุปกรณ์ที่รองรับ StarFlash มากขึ้น และได้ประสบการณ์ไร้สายที่ดีขึ้นตามไปด้วย

\begin{tcbpractice}{ลองไปเช็กดู: มือถือของคุณรองรับฟังก์ชันเหล่านี้ไหม}
\begin{enumerate}
  \item : เข้าไปในการตั้งค่าแล้วค้นหาคำว่า "ดาวเทียม" ดูว่ามีตัวเลือกเกี่ยวกับการสื่อสารผ่านดาวเทียมไหม
  \item : ค้นหา "การตรวจจับอุบัติเหตุรถชน" หรือ "SOS ฉุกเฉิน" เช็กว่าเปิดใช้งานแล้วหรือยัง
  \item : ถ้ามือถือของคุณเป็น Huawei ลองดูว่ารองรับการเชื่อมต่ออุปกรณ์ StarFlash ไหม
  \item : ศึกษาวิธีใช้ฟังก์ชันเหล่านี้ไว้ล่วงหน้า เผื่อวันจำเป็นอาจได้ช่วยชีวิต
\end{enumerate}

\end{tcbpractice}
